\documentclass[11pt]{article}

\usepackage[margin=1in]{geometry}
\usepackage{amsmath, amssymb, amsthm, bm, mathtools}
\usepackage{enumitem}
\usepackage{hyperref}

\newcommand{\E}{\mathbb{E}}
\newcommand{\Var}{\mathrm{Var}}
\newcommand{\Cov}{\mathrm{Cov}}
\newcommand{\R}{\mathbb{R}}
\newcommand{\iid}{\stackrel{\mathrm{iid}}{\sim}}
\newcommand{\Unif}{\mathrm{Unif}}
\newcommand{\Bern}{\mathrm{Bern}}

\title{Pset 1 -- Solutions}
\author{matt krasnow}
\date{February 3, 2026}

\begin{document}
\maketitle

\section*{Exercise 1 (Markov Chains)}

Let the state space be the set of cities $\{1,2,3,4,5,6,7\}$. From the figure, the (undirected) neighbor sets are
\[
\begin{aligned}
&N(1)=\{2,3,6\},\quad N(2)=\{1,4,7\},\quad N(3)=\{1,5\},\\
&N(4)=\{2,5,6,7\},\quad N(5)=\{3,4,6\},\quad N(6)=\{1,4,5\},\quad N(7)=\{2,4\}.
\end{aligned}
\]
Hence the degrees are
\[
\deg(1)=3,\ \deg(2)=3,\ \deg(3)=2,\ \deg(4)=4,\ \deg(5)=3,\ \deg(6)=3,\ \deg(7)=2,
\]
and $\sum_{j=1}^7 \deg(j)=20$.

\begin{enumerate}[label=(\alph*)]
\item \textbf{Prove that $(X_j)$ is a Markov chain.}

By construction: given $X_j=i$, the next city $X_{j+1}$ is chosen uniformly from the neighbors $N(i)$, and does not depend on earlier history.

Formally, for any $j\ge 0$ and any states $i,k$,
\[
\mathbb P(X_{j+1}=k \mid X_j=i, X_{j-1},\dots,X_0)
=
\mathbb P(X_{j+1}=k\mid X_j=i)
=
\begin{cases}
\frac{1}{\deg(i)}, & k\in N(i),\\[4pt]
0, & k\notin N(i).
\end{cases}
\]
This is the Markov property, so $(X_j)$ is a Markov chain.

\item \textbf{What is the transition matrix $P$?}

Index states in the order $(1,2,3,4,5,6,7)$. Then $P_{ij}=\mathbb P(X_{t+1}=j\mid X_t=i)$ equals $1/\deg(i)$ if $j\in N(i)$, else $0$. Thus
\[
P=
\begin{pmatrix}
0 & \tfrac13 & \tfrac13 & 0 & 0 & \tfrac13 & 0\\
\tfrac13 & 0 & 0 & \tfrac13 & 0 & 0 & \tfrac13\\
\tfrac12 & 0 & 0 & 0 & \tfrac12 & 0 & 0\\
0 & \tfrac14 & 0 & 0 & \tfrac14 & \tfrac14 & \tfrac14\\
0 & 0 & \tfrac13 & \tfrac13 & 0 & \tfrac13 & 0\\
\tfrac13 & 0 & 0 & \tfrac13 & \tfrac13 & 0 & 0\\
0 & \tfrac12 & 0 & \tfrac12 & 0 & 0 & 0
\end{pmatrix}.
\]

\item \textbf{Prove that the stationary distribution is $\pi(j)=\deg(j)/20$.}

Define
\[
\pi(j)=\frac{\deg(j)}{\sum_{\ell=1}^7 \deg(\ell)}=\frac{\deg(j)}{20}.
\]
We verify $\pi P=\pi$. For any fixed $j$,
\[
(\pi P)(j)=\sum_{i=1}^7 \pi(i)P_{ij}
= \sum_{i\in N(j)} \pi(i)\frac{1}{\deg(i)}
= \sum_{i\in N(j)} \frac{\deg(i)}{20}\cdot \frac{1}{\deg(i)}
= \sum_{i\in N(j)} \frac{1}{20}
= \frac{\deg(j)}{20}
=\pi(j).
\]
So $\pi$ is stationary.

\item \textbf{Prove that the chain is reversible: $\pi(i)P_{ij}=\pi(j)P_{ji}$.}

If $i$ and $j$ are not neighbors, then $P_{ij}=P_{ji}=0$, so the equality holds.

If $i\sim j$ (there is an edge), then
\[
\pi(i)P_{ij}=\frac{\deg(i)}{20}\cdot \frac{1}{\deg(i)}=\frac{1}{20}
=\frac{\deg(j)}{20}\cdot \frac{1}{\deg(j)}=\pi(j)P_{ji}.
\]
Hence the chain is reversible.

\item \textbf{Lazy walk: stay with probability $\alpha\in(0,1)$. Does $\pi$ change?}

The new transition matrix is
\[
P^{(\alpha)}=\alpha I + (1-\alpha)P.
\]
If $\pi$ is stationary for $P$, i.e.\ $\pi P=\pi$, then
\[
\pi P^{(\alpha)}
=\pi\bigl(\alpha I + (1-\alpha)P\bigr)
=\alpha \pi + (1-\alpha)\pi P
=\alpha \pi + (1-\alpha)\pi
=\pi.
\]
So the stationary distribution does \emph{not} change: it remains $\pi(j)=\deg(j)/20$.
\end{enumerate}

\section*{Exercise 2 (Linear Algebra)}

From the figure, the undirected edges are
\[
\{(1,2),(1,3),(2,3),(2,4),(3,5),(4,5),(4,6),(5,6),(4,7),(6,7)\}.
\]
Degrees:
\[
\deg(1)=2,\ \deg(2)=3,\ \deg(3)=3,\ \deg(4)=4,\ \deg(5)=3,\ \deg(6)=3,\ \deg(7)=2.
\]

\begin{enumerate}[label=(\alph*)]
\item \textbf{Adjacency matrix $A$ and Laplacian $L=D-A$.}

With vertex order $(1,2,3,4,5,6,7)$,
\[
A=
\begin{pmatrix}
0&1&1&0&0&0&0\\
1&0&1&1&0&0&0\\
1&1&0&0&1&0&0\\
0&1&0&0&1&1&1\\
0&0&1&1&0&1&0\\
0&0&0&1&1&0&1\\
0&0&0&1&0&1&0
\end{pmatrix},
\qquad
D=\mathrm{diag}(2,3,3,4,3,3,2).
\]
Thus
\[
L=D-A=
\begin{pmatrix}
2&-1&-1&0&0&0&0\\
-1&3&-1&-1&0&0&0\\
-1&-1&3&0&-1&0&0\\
0&-1&0&4&-1&-1&-1\\
0&0&-1&-1&3&-1&0\\
0&0&0&-1&-1&3&-1\\
0&0&0&-1&0&-1&2
\end{pmatrix}.
\]

\item \textbf{Compute the eigenvalues of $L$. Remark that $\bm 1$ is an eigenvector.}

Numerically (e.g.\ using Python for the symmetric matrix $L$), the eigenvalues are approximately
\[
\lambda(L)\approx
\{0,\ 0.814349,\ 2.328009,\ 3.313908,\ 3.598089,\ 4.457530,\ 5.488115\}.
\]
Also,
\[
L\bm 1 = 0\quad\text{where }\bm 1=(1,1,1,1,1,1,1)^\top,
\]
because each row of $L$ sums to $0$. Hence $\bm 1$ is an eigenvector with eigenvalue $0$.

\item \textbf{For a general undirected graph, prove that for any $x\in\R^n$,
\(
x^\top L x = \sum_{(i,j)\in E}(x_i-x_j)^2.
\)
}

For an undirected graph,
\[
x^\top L x = x^\top(D-A)x
= \sum_{i}\deg(i)x_i^2 - \sum_{i,j}A_{ij}x_i x_j.
\]
Rewrite the first term as $\sum_{i,j}A_{ij}x_i^2$ (since $\deg(i)=\sum_j A_{ij}$):
\[
x^\top L x
= \sum_{i,j}A_{ij}x_i^2 - \sum_{i,j}A_{ij}x_i x_j
= \frac12\sum_{i,j}A_{ij}\bigl(x_i^2+x_j^2-2x_i x_j\bigr)
= \frac12\sum_{i,j}A_{ij}(x_i-x_j)^2.
\]
Each undirected edge $(i,j)\in E$ appears twice in $\sum_{i,j}$, so
\[
x^\top L x = \sum_{(i,j)\in E}(x_i-x_j)^2.
\]

\item \textbf{Use this identity to prove that $L$ is positive semidefinite and that the smallest eigenvalue is always $\lambda_1=0$.}

From (c),
\[
x^\top L x = \sum_{(i,j)\in E}(x_i-x_j)^2 \ge 0\quad\text{for all }x,
\]
so $L\succeq 0$ (positive semidefinite), hence all eigenvalues are $\ge 0$.
But $L\bm 1=0$, so $0$ is an eigenvalue. Therefore the smallest eigenvalue is
\[
\lambda_1=0.
\]

\item \textbf{Let $\lambda_2$ be the second-smallest eigenvalue. Prove that if $G$ is disconnected, then $\lambda_2=0$.}

Assume $G$ has at least two connected components, and let $C$ be one component. Let $\bm 1_C\in\R^n$ be the indicator vector:
\[
(\bm 1_C)_i=
\begin{cases}
1,& i\in C,\\
0,& i\notin C.
\end{cases}
\]
Because there are no edges between $C$ and $C^c$, every edge $(i,j)\in E$ has both endpoints in the same component, hence $(\bm 1_C)_i-(\bm 1_C)_j=0$. By (c),
\[
\bm 1_C^\top L \bm 1_C = \sum_{(i,j)\in E}\bigl((\bm 1_C)_i-(\bm 1_C)_j\bigr)^2=0,
\]
so $\bm 1_C\in \ker(L)$. Also $\bm 1_C$ is not a multiple of $\bm 1$, so $\dim \ker(L)\ge 2$, meaning eigenvalue $0$ has multiplicity at least $2$. Therefore the second-smallest eigenvalue is
\[
\lambda_2=0.
\]
\end{enumerate}

\section*{Exercise 3 (Statistical Inference)}

Let $X_1,\dots,X_n$ be i.i.d.\ with density
\[
f(x)=\frac{1}{2\sigma}e^{-|x|/\sigma},\qquad \sigma>0.
\]

\begin{enumerate}[label=(\alph*)]
\item \textbf{Write out the log-likelihood function $\ell(X_1,\dots,X_n)$.}

The likelihood is
\[
L(\sigma)=\prod_{i=1}^n \frac{1}{2\sigma}\exp\!\left(-\frac{|X_i|}{\sigma}\right).
\]
Thus the log-likelihood is
\[
\ell(\sigma)=\log L(\sigma)
= \sum_{i=1}^n\left[-\log(2\sigma)-\frac{|X_i|}{\sigma}\right]
= -n\log(2\sigma)-\frac{1}{\sigma}\sum_{i=1}^n |X_i|.
\]

\item \textbf{What is the MLE of $\sigma$?}

Differentiate for $\sigma>0$:
\[
\ell'(\sigma)= -\frac{n}{\sigma}+\frac{1}{\sigma^2}\sum_{i=1}^n |X_i|.
\]
Set $\ell'(\sigma)=0$:
\[
-\frac{n}{\sigma}+\frac{1}{\sigma^2}\sum_{i=1}^n |X_i|=0
\quad\Longleftrightarrow\quad
-n\sigma+\sum_{i=1}^n|X_i|=0
\quad\Longleftrightarrow\quad
\hat\sigma=\frac1n\sum_{i=1}^n |X_i|.
\]
Second derivative:
\[
\ell''(\sigma)=\frac{n}{\sigma^2}-\frac{2}{\sigma^3}\sum_{i=1}^n |X_i|,
\]
and at $\sigma=\hat\sigma$, $\ell''(\hat\sigma)=\frac{n}{\hat\sigma^2}-\frac{2n\hat\sigma}{\hat\sigma^3}=-\frac{n}{\hat\sigma^2}<0$, so $\hat\sigma$ is indeed a maximizer.

\item \textbf{Is the MLE consistent? Prove your answer.}

Compute $\E|X_1|$. By symmetry,
\[
\E|X_1|
=2\int_0^\infty x\cdot \frac{1}{2\sigma}e^{-x/\sigma}\,dx
=\frac{1}{\sigma}\int_0^\infty x e^{-x/\sigma}\,dx.
\]
Let $y=x/\sigma$ so $x=\sigma y$, $dx=\sigma dy$:
\[
\E|X_1|
=\frac{1}{\sigma}\int_0^\infty (\sigma y)e^{-y}\,\sigma dy
=\sigma\int_0^\infty y e^{-y}\,dy
=\sigma\cdot 1
=\sigma,
\]
since $\int_0^\infty y e^{-y}\,dy=1$.

Now
\[
\hat\sigma=\frac{1}{n}\sum_{i=1}^n |X_i|.
\]
By the strong law of large numbers,
\[
\hat\sigma \xrightarrow{a.s.} \E|X_1|=\sigma,
\]
so $\hat\sigma$ is consistent.
\end{enumerate}

\section*{Exercise 4 (Probability Theory)}

Assume the standard (undirected) graphon sampling: draw $U_1,\dots,U_n \iid \Unif(0,1)$, and (conditionally on the $U$'s) the edges $\{a_{ij}:1\le i<j\le n\}$ are independent with
\[
a_{ij}\mid (U_i,U_j)\sim \Bern(W(U_i,U_j)),\qquad a_{ji}=a_{ij},\qquad a_{ii}=0.
\]
Define degrees $d_i=\sum_{j\ne i} a_{ij}$.

Introduce the graphon degree function
\[
g(u):=\int_0^1 W(u,v)\,dv,
\]
and the global edge density
\[
\mu := \int_0^1\int_0^1 W(u,v)\,du\,dv = \E[W(U_1,U_2)].
\]

\begin{enumerate}[label=(\alph*)]
\item \textbf{Compute $\E[d_i\mid U_i]$ and $\E[d_i]$.}

For $j\ne i$,
\[
\E[a_{ij}\mid U_i]
=\E\!\left[\E[a_{ij}\mid U_i,U_j]\mid U_i\right]
=\E[W(U_i,U_j)\mid U_i]
=\int_0^1 W(U_i,v)\,dv
=g(U_i).
\]
Therefore
\[
\E[d_i\mid U_i]=(n-1)\,g(U_i).
\]
Taking expectation again,
\[
\E[d_i]=(n-1)\,\E[g(U_i)]=(n-1)\mu.
\]

\item \textbf{Compute $\Var(d_i\mid U_i)$ and $\Var(d_i)$.}

Given $U_i$, the random variables $\{a_{ij}:j\ne i\}$ are independent and each satisfies
\[
\mathbb P(a_{ij}=1\mid U_i)=g(U_i),
\]
so $a_{ij}\mid U_i \sim \Bern(g(U_i))$. Hence
\[
\Var(a_{ij}\mid U_i)=g(U_i)\bigl(1-g(U_i)\bigr),
\]
and summing independent terms,
\[
\Var(d_i\mid U_i)=(n-1)\,g(U_i)\bigl(1-g(U_i)\bigr).
\]

For the unconditional variance, use the law of total variance:
\[
\Var(d_i)=\E[\Var(d_i\mid U_i)] + \Var(\E[d_i\mid U_i]).
\]
Thus
\[
\Var(d_i)
=(n-1)\,\E\!\big[g(U)(1-g(U))\big] + (n-1)^2\,\Var(g(U)).
\]

\item \textbf{Compute $\Cov(d_i,d_j)$ for $i\ne j$.}

Write
\[
d_i = a_{ij} + \sum_{k\ne i,j} a_{ik},
\qquad
d_j = a_{ij} + \sum_{k\ne i,j} a_{jk}.
\]
Let
\[
\gamma := \E\!\left[g(U)^2\right]
= \E\!\left[W(U_1,U_2)\,W(U_1,U_3)\right]
\quad (\text{with }U_1,U_2,U_3\ iid).
\]
Then
\[
\E[a_{ij}]=\mu,\quad \Var(a_{ij})=\mu(1-\mu),
\]
and for any $k\ne i,j$,
\[
\Cov(a_{ij},a_{ik})=\gamma-\mu^2,\quad
\Cov(a_{ij},a_{jk})=\gamma-\mu^2,\quad
\Cov(a_{ik},a_{jk})=\gamma-\mu^2.
\]
There are $(n-2)$ choices of $k$, so
\[
\Cov(d_i,d_j)=\mu(1-\mu)\;+\;3(n-2)\,(\gamma-\mu^2).
\]

\item \textbf{Limits of $d_1/n$ and $d_2/n$. Do they converge to the same quantity? CLT?}

Condition on $U_1$. Since $d_1\mid U_1 \sim \mathrm{Bin}(n-1,g(U_1))$,
\[
\frac{d_1}{n-1}\xrightarrow{a.s.} g(U_1),
\qquad\text{hence}\qquad
\frac{d_1}{n}\xrightarrow{a.s.} g(U_1).
\]
Similarly,
\[
\frac{d_2}{n}\xrightarrow{a.s.} g(U_2).
\]
In general, they do \emph{not} converge to the same quantity: $g(U_1)$ and $g(U_2)$ are i.i.d.\ and are almost surely equal only if $g$ is (a.e.) constant.

\medskip
\noindent
\textbf{CLT (conditional).} For fixed $u$, if $d_1\mid U_1=u\sim \mathrm{Bin}(n-1,g(u))$, then
\[
\frac{d_1-(n-1)g(u)}{\sqrt{(n-1)g(u)(1-g(u))}}
\xrightarrow{d} \mathcal N(0,1).
\]
Equivalently, conditioning on $U_1$,
\[
\sqrt{n}\left(\frac{d_1}{n}-g(U_1)\right)\ \Big|\ U_1
\ \xrightarrow{d}\ \mathcal N\!\left(0,\ g(U_1)(1-g(U_1))\right).
\]
Unconditionally this yields a mixture of normals (variance depends on the random $U_1$).

\item \textbf{Limit of $\dfrac{\sum_{i\le n} d_i}{n^2}$.}

In an undirected graph,
\[
\sum_{i=1}^n d_i = 2\sum_{1\le i<j\le n} a_{ij}.
\]
Thus
\[
\frac{1}{n^2}\sum_{i=1}^n d_i
= \frac{2}{n^2}\sum_{i<j} a_{ij}.
\]
Decompose $a_{ij}=(a_{ij}-W(U_i,U_j))+W(U_i,U_j)$. The first part averages to $0$ and the second part is a $U$-statistic that converges almost surely to $\mu=\iint W$. Therefore
\[
\frac{1}{n^2}\sum_{i=1}^n d_i \xrightarrow{a.s.} \iint_{[0,1]^2} W(u,v)\,du\,dv.
\]
\end{enumerate}

\end{document}
